\title{Parameter-Efficient Orthogonal Finetuning via Factorization}

\begin{abstract}
The prevalence of large foundation models continues to grow, yet training these models from the ground up remains infeasible due to computational constraints. This necessitates strategies for adapting these high-capacity models to new tasks in a resource-efficient manner. In this work, we investigate Orthogonal Finetuning (OFT), a methodically grounded paradigm for adapting models to downstream applications. Although OFT has been shown to generalize well, it typically requires many tunable parameters owing to the intrinsically high dimensionality of orthogonal matrices. To mitigate this, we analyze OFT through the lens of information transfer and pinpoint several requirements for enhancing parameter efficiency. Drawing inspiration from the Cooley-Tukey fast Fourier transform algorithm's ability to transmit information with reduced complexity, we develop an efficient orthogonal parameterization leveraging butterfly structures. Incorporating this technique into OFT yields a new, parameter-efficient finetuning approach, dubbed Orthogonal Butterfly~(BOFT). BOFT serves as a generalization of OFT, encompassing it as a particular instance, and provides a broader orthogonal finetuning framework. Comprehensive empirical evaluations are conducted on the adaptation of large vision transformers, language models, and text-to-image diffusion models to a diverse array of vision and language benchmarks.
\end{abstract}

\section{Introduction}

Models like ChatGPT~\cite{brown2020language,chen2021evaluating} and Stable Diffusion~\cite{rombach2022high} exemplify the extraordinary capacity of large foundation models to generalize across tasks. However, the substantial improvements in their effectiveness have been accompanied by explosive growth in parameter counts—

(\emph{e.g.}, GPT-3 comprises roughly 175B parameters)—making training these architectures from the beginning increasingly impractical for most researchers. Consequently, further advancement in the field hinges on the development of methods that allow for effective adaptation of such foundation models without necessitating full retraining. Efficient adaptation for downstream tasks has thus become a primary concern. The main strategies for achieving this goal are: \emph{model finetuning}~\cite{hulora2022,zaken2022bitfit,guo2020parameter,raffel2020exploring,qiu2023controlling,zhang2022adaptive,chavan2023one}, where only select parameters within the fixed architecture are updated; \emph{adapter tuning}~\cite{rebuffi2017learning,houlsby2019parameter,pfeiffer2020adapterfusion,lin2020exploring,he2021towards}, which augments the existing model with supplementary trainable modules; and \emph{prompt tuning}~\cite{li2021prefix,lester2021power}, wherein trainable prefix tokens are attached to the input. Of these methods, model finetuning stands out due to its straightforwardness and efficiency, notably adding no computational overhead during inference.

The core idea underlying model finetuning is to maintain the similarity between the pretrained and finetuned models, according to specific metrics, so as to retain the knowledge acquired during pretraining. A typical strategy in current finetuning practices involves selecting a low learning rate, which implicitly restricts the difference between the original and the updated model parameters. Recognizing that weight matrices possess inherent structure, a more rigorous method seeks to conserve the relational characteristics of these matrices, particularly the pairwise angles among neurons. This consideration has inspired the development of a new finetuning paradigm, termed Orthogonal Finetuning~(OFT)~\cite{qiu2023controlling}. In OFT, all neurons within a single layer undergo the same orthogonal transformation, guaranteeing the preservation of pairwise angular relationships throughout the finetuning stage. While OFT has shown substantial improvements in generalizability and convergence, particularly for text-to-image diffusion model finetuning~\cite{qiu2023controlling}, the approach entails a substantial number of learnable parameters due to the large size of orthogonal matrices. To mitigate this, OFT employs a block-diagonal configuration that diminishes the parameter count. Nevertheless, this approach introduces certain limitations: the fixed sparsity structure of the orthogonal matrices leads to independent application of orthogonal transformations within each block. Although such block-diagonal sparsity enhances parameter efficiency, it may inadvertently induce unwanted inductive biases (for instance, the restriction reduces the matrix’s expressivity and its ability to emulate standard linear operations). Furthermore, determining an optimal sparsity structure remains an unresolved issue.

A central challenge here is to construct a dense orthogonal matrix while minimizing the number of trainable parameters. At first, this seems difficult, as a dense orthogonal matrix in $d$ dimensions typically demands $\mathcal{O}(d^2)$ parameters. However, we propose an alternative approach: instead of representing the matrix directly, we build it as a product of several factorized sparse matrices. This strategy rests on the observation that computation can be increased to reduce the number of parameters stored—by composing the orthogonal matrix through sequential multiplications of sparse matrices, we offload the parameter overhead to repeated computation at each training step. To further conceptualize this, we frame dense orthogonal matrix construction as an information dissemination problem, where generating the dense matrix via multiplication of sparse components mirrors information transfer within a grid-like graph. Viewing the problem this way uncovers a variety of possible sparse matrix factorizations that constrain parameter count but retain sufficient capacity to create dense orthogonal matrices.

To optimize parameter usage in this information transmission paradigm, we are inspired by the butterfly network used in the Cooley-Tukey fast Fourier transform algorithm~\cite{cooley1965algorithm}, known for its efficient communication between nodes~\cite{klasing1994broadcasting}. The butterfly network of the Cooley-Tukey algorithm directly informs a type of sparse matrix decomposition, referred to as butterfly factorization. Employing butterfly factorization allows us to express a $d$-dimensional dense matrix as a sequence of $\mathcal{O}(\log d)$ sparse matrices, each containing just $\mathcal{O}(d)$ nonzero elements. When these constituent sparse matrices are orthogonal, we obtain a highly compact orthogonal parameterization, utilizing only $\mathcal{O}(d\log d)$ parameters—a marked improvement over conventional approaches. Building on this compact representation, we introduce a new, parameter-efficient finetuning strategy, Orthogonal Butterfly~(BOFT), which encompasses OFT as a special instance and establishes a broad orthogonal finetuning framework. Both the block-diagonal structure and the butterfly architecture share a core feature—sparsity—which underlies their ability to decrease the parameter footprint of orthogonal matrices. This perspective offers an interesting juxtaposition with the low-rank matrices used in LoRA~\cite{hulora2022}; notably, low-rank and sparse matrices are two principal structured matrix categories~\cite{candes2011robust} that promote parameter efficiency through different structural constraints.

In contrast to OFT's use of a block-diagonal structure—which balances model expressivity and regularity—BOFT adopts a butterfly structure that enables more gradual transitions between matrices in the full orthogonal group (allowing for higher expressivity) and the identity matrix (favoring regularity). This approach helps restrict the hypothesis space within the orthogonal group for downstream applications. Given the butterfly structure’s prominent role in efficient linear operators like the discrete Fourier and cosine transforms, we posit that our approach to parameter efficiency introduces an implicit inductive bias, thus enhancing generalization and reducing overfitting risk. Our reasoning is grounded in the observation that while the butterfly structure can reconstruct various well-known linear transforms, the block-diagonal structure utilized by OFT cannot. The principal contributions of our work are outlined as follows:

\begin{itemize}[leftmargin=*,nosep]
    \item We address parameter efficiency in orthogonal finetuning by developing a new information transmission perspective. This paradigm reframes the creation of compact orthogonal matrices as an information flow problem on a grid-like graph.
    \item Drawing from the butterfly patterns seen in the Cooley-Tukey algorithm, we introduce Orthogonal Butterfly—a highly parameter-efficient orthogonal finetuning strategy—within our information transmission framework.
    \item We present several theoretical arguments for how BOFT drastically reduces the parameter count while preserving both expressivity and training stability. The underlying matrix factorization further enables a compelling scheme for weight interpolation (as illustrated in Figure~\ref{fig:boft_interpolation}).
    \item For the first time, we extend the use of orthogonal finetuning beyond controllable text-to-image generation~\cite{qiu2023controlling}, establishing its versatility as a general approach for model adaptation. In pursuit of this goal, we deploy BOFT across diverse domains, including computer vision and natural language processing. Our method consistently surpasses the current state-of-the-art in terms of both parameter efficiency and generalization, as evidenced by our experimental results.
\end{itemize}

\section{Related Work}

\textbf{Parameter-efficient finetuning (PEFT)}. As foundation models (\emph{e.g.}, \cite{brown2020language,radford2021learning,rombach2022high,kirillov2023segment}) continue to scale in size and capability, efficiently adapting these large models to downstream tasks without modifying all parameters has attracted substantial attention~\cite{treviso2023efficient,ding2023parameter,lialin2023scaling}. Numerous PEFT techniques have been proposed~\cite{houlsby2019parameter,aghajanyan2020intrinsic,hulora2022,edalati2022krona,wang2022adamix,gheini2021cross,zaken2022bitfit,guo2020parameter,sung2021training,ansell2022composable,lester2021power,li2021prefix,vu2022spot,he2021towards,mao2021unipelt,karimi2021compacter,liu2022few,sung2022lst,chen2022parameter,jia2022visual,chen2022adaptformer,zhang2022neural,jie2023fact,lian2022scaling,luo2023towards}, with reparameterization-based strategies~\cite{aghajanyan2020intrinsic,hulora2022,edalati2022krona} being particularly pertinent to our approach. LoRA~\cite{hulora2022}, for example, finetunes pretrained weight matrices by incorporating the product of two low-rank matrices, yielding strong results in natural language domains. However, LoRA employs a fixed-rank parameterization for all matrices; as a result, subsequent methods~\cite{zhang2022adaptive,valipour2022dylora,zhang2023increlora} have introduced adaptive rank adjustments per layer to optimize parameter allocation under a given budget. This straightforward low-rank approach to weight modification has become very widespread~\cite{dettmers2023qlora,zi2023delta,chavan2023one}. Motivated by findings linking hyperspherical energy to generalization performance~\cite{liu2018learning,liu2021orthogonal}, \cite{qiu2023controlling} introduced orthogonal finetuning, a distinct and effective strategy for adapting text-to-image diffusion models. In this method, an orthogonal transformation is learned for neurons within a single layer, leading to improved generalization and more stable training relative to LoRA. Despite these advantages, OFT generally introduces more tunable parameters compared to LoRA, making the pursuit of more parameter-efficient variants highly valuable. Furthermore, the broader applicability of OFT beyond its original scope—namely, controlling text-to-image diffusion models—remains unexplored. BOFT addresses these concerns by employing butterfly factorization to enhance the parameter efficiency of OFT. This advancement enables us to empirically investigate orthogonal finetuning’s utility across a variety of adaptation settings.

\textbf{Butterfly structures}. The radix-2 Cooley-Tukey procedure successively decomposes the $N$-point discrete Fourier transform into a pair of $\frac{N}{2}$-point transforms at each step, naturally inducing what is known as a butterfly structure—expressible as a sequence (product) of sparse matrices, collectively referred to as a butterfly matrix. Such matrices have been leveraged for parameterizing orthogonal transforms to circumvent the need for pivoting during Gaussian elimination and to enhance computational efficiency~\cite{parker1995random}, for improving the training stability of recurrent neural networks~\cite{jing2017tunable}, as well as in kernel approximation contexts~\cite{munkhoeva2018quadrature}. The works~\cite{dao2019learning,dao2019kaleidoscope} explore fast linear operators built with butterfly parameterizations, while~\cite{chen2022pixelated,dao2022monarch} exploit them for the accelerated training of neural networks. Additionally, butterfly architectures facilitate efficient matrix-vector multiplication~\cite{michielssen1996multilevel,o2010algorithm}, compact matrix approximations~\cite{li2015butterfly}, and high-throughput network communication~\cite{klasing1994broadcasting,li2003linear,rajasingh2016transmission}. Diverging from prior applications, our investigation centers on utilizing butterfly structures specifically to improve the parameter efficiency of OFT.

\section{An Information Transmission View on Orthogonal Finetuning}

We begin by reviewing the foundational concepts of OFT. In order to fine-tune a pretrained weight matrix, OFT re-expresses the updated weight matrix as the product of a learnable orthogonal matrix and the original, frozen pretrained weight matrix. Unlike LoRA, which employs an additive low-rank update to the weights, OFT adopts a multiplicative strategy via an orthogonal matrix. LoRA attains parameter efficiency through the low-rank structure of its updates, whereas the classic OFT achieves this by using a block-diagonal form for its orthogonal matrices~\cite{qiu2023controlling}; the smaller the diagonal block size, the greater the parameter efficiency obtained by OFT. Figure~\ref{fig:comp_lora} provides a clear visualization of this comparison. The primary rationale for leveraging an orthogonal transformation in fine-tuning is to preserve the pairwise angular relationships between neurons~\cite{liu2018learning,liu2021orthogonal,qiu2023controlling}, thereby maintaining the semantic knowledge embedded during pretraining. Specifically, OFT seeks to learn an orthogonal matrix $\bm{R}\in\mathbb{R}^{d\times d}$ for a pretrained linear transformation $\bm{W}^0\in\mathbb{R}^{d\times n}$, and modifies the computation in the forward pass from $\bm{z}=(\bm{W}^0)^\top\bm{x}$ to $\bm{z}=(\bm{R}\bm{W}^0)^\top\bm{x}$, where the vectors $\bm{x}\in\mathbb{R}^{d}$ and $\bm{z}\in\mathbb{R}^{n}$ serve as input and output, respectively. To guarantee that fine-tuning begins from the pretrained model’s original behavior, $\bm{R}$ is initially set to the identity matrix. Maintaining $\bm{R}$ as orthogonal during optimization follows the Cayley parameterization strategy, as described in \cite{qiu2023controlling,liu2021orthogonal}; specifically, $\bm{R}=(\bm{I}+\bm{Q})(\bm{I}-\bm{Q})^{-1}$, where $\bm{Q}$ is skew-symmetric, satisfying $\bm{Q}=-\bm{Q}^\top$. For the sake of parameter efficiency, a block-diagonal parameterization is imposed on $\bm{R}$, expressed as $\text{diag}(\bm{R}_1,\bm{R}_2,\cdots,\bm{R}_r)$ with each $\bm{R}_i\in\mathcal{R}^{b\times b}$ being an orthogonal block, where $br=d$. While this block-diagonal construction reduces parameter count, it is based on the simplifying assumption that neuron dimensions—corresponding to columns in $\bm{W}^0$—are partitioned into $r$ groups, each transformed independently by a separate orthogonal matrix. Empirical results have shown the effectiveness of this design, but there is little justification for splitting a neuron’s dimensions into groups merely according to their indices. This limitation highlights the appeal of dense orthogonal matrices. Hence, a compelling research direction is posed: \emph{Is it feasible to realize a dense orthogonal matrix that retains parameter efficiency?}

In response to this problem, we introduce a method in which a dense orthogonal matrix is constructed as the product of several sparse orthogonal matrices. To systematically analyze the sparsity structures that arise in orthogonal matrix factorization, we reinterpret the construction of a dense orthogonal matrix within OFT as an instance of information transmission. Concretely, when a dense matrix $\bm{R}\in\mathbb{R}^{d\times d}$ is expressed as a sequence of $m$ square matrices, namely $\thickmuskip=2mu \medmuskip=2mu \bm{R}=\bm{B}_m\bm{B}_{m-1}\cdots\bm{B}_1$, this can be visualized as information propagation across a grid consisting of $d\times (m+1)$ nodes, as depicted in Figure~\ref{fig:info_trans}. The intuition for this perspective arises from treating a $d$-dimensional dense matrix as representing full connectivity between two sets of $d$ nodes. Specifically, if the entry $R_{ij}$ in $\bm{R}$ is zero, it signifies the impossibility of transmitting information from node $j$ to node $i$, whereas a nonzero value permits such transmission. As a result, expressing $\bm{R}$ through a product of matrices $\bm{B}_m\bm{B}_{m-1}\cdots\bm{B}_1$ naturally equates to implementing a series of information exchanges, governed by the connectivity graphs determined by each $\bm{B}_i$. The order of these exchanges follows $\bm{B}_1$ initially and culminates with $\bm{B}_n$.

To exemplify, Figure~\ref{fig:info_trans} visualizes the case $\bm{R}=\bm{B}_5\bm{B}_4\bm{B}_3\bm{B}_2\bm{B}_1$, illustrating both the respective sparsity structures and the associated induced graph. The depicted graph results from the stepwise unfolding of the matrix multiplications. Within this structure, every matrix $\bm{B}_i$ defines connectivity from the $i$-th set of nodes to the next, i.e., the $(i+1)$-th layer. Concretely, the element at $(j_1, j_2)$ in $\bm{B}_i$ indicates the presence or absence (through nonzero or zero) of a directed edge from node $j_2$ at level $i$ to node $j_1$ at level $(i+1)$. Ensuring the product $\bm{B}_5\bm{B}_4\bm{B}_3\bm{B}_2\bm{B}_1$ yields a dense matrix requires that each node in the initial level be capable of sending information to every node in the sixth level. By focusing on the truncated factorization $\bm{R}=\bm{B}_4\bm{B}_3\bm{B}_2\bm{B}_1$ (linking the nodes in levels one and five), it is evident, for example, that node 1 is unable to transmit to node 3; this corresponds to a zero at position $(3,1)$ in the compound sparsity pattern. The same relationship between the graph structure and matrix sparsity extends to shorter products, such as $\bm{R}=\bm{B}_3\bm{B}_2\bm{B}_1$ and $\bm{R}=\bm{B}_2\bm{B}_1$. More generally, achieving a dense matrix $\bm{R}\in\mathbb{R}^{d\times d}$ requires that the sequence $\bm{B}_m,\ldots,\bm{B}_1$ collectively defines a system of directed edges on a $d\times (m+1)$ grid, connecting nodes exclusively between successive levels (that is, between adjacent columns), so that any node in the initial level can ultimately communicate with every node at the final, $(m+1)$-st, level.

The information transmission perspective provides valuable insights into the sparsity structure of factorization matrices within OFT. As depicted in Figure~\ref{fig:blk_diag}, the matrix $\bm{R}$ in standard OFT exhibits a block-diagonal layout. While this configuration reduces the quantity of parameters requiring optimization, it cannot yield a densely connected $\bm{R}$. The main objective is to synthesize a dense orthogonal matrix by composing $m$ sparse orthogonal matrices, minimizing the number of trainable parameters involved. Through the lens of information transmission, two fundamental requirements emerge: \emph{(i)} \textbf{dense connectivity}, ensuring each node in the initial layer connects, through some path, to every node in the final layer; and \emph{(ii)} \textbf{minimum free edges}, such that, under the orthogonality constraint, the overall number of (trainable) edges remains minimal. It is important to observe that the orthogonality constraint imposes a specific structure on the edges linking adjacent layers. Notably, for each $\bm{B}_i$ to maintain full rank—a prerequisite for orthogonality—there must be $d$ connections forming a bijection between nodes in the $i$-th and $(i=1)$-th levels, guaranteeing at least $d$ edges between two neighboring layers (see the example in Figure~\ref{fig:info_trans}, where $d=4$). These $d$ connections, which uphold orthogonality and are not subject to training (for instance, a $d\times d$ orthogonal matrix with $d$ nonzero entries, each constrained to be $\pm 1$), are not counted as trainable edges. Since orthogonal matrices require fewer independent parameters compared to generic dense matrices, enforcing orthogonality introduces dependencies among edge connections. For instance, in a $2\times 2$ orthogonal matrix, placing a zero in position $(1,1)$ necessitates another zero in $(2,2)$—indicating that the removal of one edge can compel the removal of another. Consequently, depending on the configuration of the edge set, orthogonality might either introduce or eliminate connections. By examining the distribution of non-zero entries in the constructed orthogonal matrix, the information transmission approach proves especially instructive in the context of OFT, where attention centers on the trainable nonzero elements of $\bm{R}$, without regard for their specific magnitudes.

A straightforward approach to densely connecting two layers requires $\mathcal{O}(d^2)$ edges, corresponding to a single dense orthogonal matrix, resulting in $d^2-d$ individual edges (for instance, when $d=4$, there are 12 such edges). Figure~\ref{fig:info_trans} depicts an example where a matrix is factorized in a feasible way, utilizing only $10$ edges—fewer than the number needed for one full dense orthogonal connection. Through this perspective, we are able to analyze the parameter efficiency of OFT via graphical representation, which also facilitates the generation of viable factorizations. We take cues from the Cooley-Tukey algorithm’s network design, particularly the butterfly graph~\cite{cooley1965algorithm}, renowned for establishing dense interconnections between $d$ sources and $d$ receivers using just $\mathcal{O}(d\log d)$ edges. To illustrate, the configuration in Figure~\ref{fig:info_trans} employs $10$ edges for dense connectivity, whereas the butterfly network achieves this with only $8$. In what follows, we detail how the butterfly structure enhances parameter efficiency.

\section{Orthogonal Parameterization by Butterfly Factorization}

Initially introduced in the Cooley-Tukey algorithm for efficient computation of the fast Fourier transform, the butterfly structure underpins rapid domain conversions—where a minor alteration in the frequency domain can propagate globally across the spatial domain. This is analogous to our information transmission setting: every node in the input layer is capable of conveying information to each node in the output layer. Additionally, the butterfly topology is well-established in the design of computer networks~\cite{solihin2015fundamentals,li2011linear}, supporting streamlined data exchange. Let $k \geq 2$ be a power of $2$; we define the \emph{butterfly factor} $\bm{B}^F(k)$ as follows:

\begin{equation}\label{eq:but_factor}
\begin{aligned}
    \bm{B}^F(k)=\begin{bmatrix}
    \text{diag}(\bm{d}_1) & \text{diag}(\bm{d}_2) \\
    \text{diag}(\bm{d}_3) & \text{diag}(\bm{d}_4)
    \end{bmatrix}\in\mathbb{R}^{k\times k},
\end{aligned}
\end{equation}

where $\bm{d}_i\in\mathbb{R}^{\frac{k}{2}}$ for all $i$ are specific vectors. Given $d=2^N$, the $d$-dimensional \emph{butterfly matrix} $\bm{B}(d)\in\mathbb{R}^{d\times d}$ is then defined recursively:

\begin{equation}
\begin{aligned}
    \!\!\bm{B}(d)=\tilde{\bm{B}}(d,d)\!\cdot\!\begin{bmatrix}
    \bm{B}_1(\frac{d}{2})\!\!\!\!\! & \bm{0} \\
    \bm{0}\!\!\!\!\! &\bm{B}_2(\frac{d}{2})
    \end{bmatrix}= \tilde{\bm{B}}(d,d)\tilde{\bm{B}}(d,\frac{d}{2})\cdots\tilde{\bm{B}}(d,2),
\end{aligned}
\end{equation}

where $\bm{B}_1(\frac{d}{2})$ and $\bm{B}_2(\frac{d}{2})$ denote two butterfly matrices of dimension $\frac{d}{2}$. We introduce the concept of the \emph{butterfly component}, expressed as $\thickmuskip=2mu \medmuskip=2mu \tilde{\bm{B}}(d,k)=\text{diag}(\bm{B}^F_1(k),\cdots,\bm{B}^F_{d/k}(k))$, a block-diagonal matrix of size $d\times d$ composed of blocks of size $k$, where each block $\bm{B}^F_i(k),\forall i$ refers to the butterfly factors given in Equation~\ref{eq:but_factor}. This structure enables us to use the butterfly matrix to parameterize orthogonal matrices. To do so, it suffices to ensure that every multiplicative element $\tilde{\bm{B}}(d,k),\forall k$ in $\bm{B}(d)$ maintains orthogonality. Starting with the block-diagonal matrix $\tilde{\bm{B}}(d,2)$, which consists of $2\times 2$ blocks, orthogonality for $\tilde{\bm{B}}(d,2)$ is easily obtained by employing the Cayley transform (or two-dimensional rotation) for each block, as described in \cite{liu2021orthogonal,qiu2023controlling}. Furthermore, the sparsity structure of any butterfly component can be interpreted as a permutation of the sparsity pattern in $\tilde{\bm{B}}(d,2)$, making the orthogonal parameterization of all butterfly components straightforward. This approach yields an efficient parameterization for orthogonal matrices by assembling multiple $2\times 2$ orthogonal blocks. Extending the butterfly matrices in line with \cite{chen2022pixelated}, we define the block butterfly component $\tilde{\bm{B}}^b(d,k)$ where each element in $\bm{d}_i,\forall i$ is a $b\times b$ matrix. To enforce orthogonality in the block butterfly component $\tilde{\bm{B}}^b(d,2)$, we parameterize each $2b\times 2b$ block to be orthogonal. The sparsity configuration of the higher-order butterfly components $\tilde{\bm{B}}^b(d,k)$ for $k>2$ results from block-wise permutations of the pattern in $\tilde{\bm{B}}^b(d,2)$ and can therefore also be made orthogonal in a similar way. Putting all the components together, the forward computation in BOFT is

\begin{equation}
    \begin{aligned}\nonumber
    \bm{z} = \big(\bm{R}(m,b) \cdot \bm{W}^0\big)^\top \bm{x},~~~~\text{s.t.}~\bigg\{\bm{R}(m,b) = \prod_{i=1}^m \tilde{\bm{B}}^b_{(d,i)}~~\&~~\underbrace{\big(\tilde{\bm{B}}^b_{(d,j)}\big)^\top\tilde{\bm{B}}^b_{(d,j)} = \tilde{\bm{B}}^b_{(d,j)}\big(\tilde{\bm{B}}^b_{(d,j)}\big)^\top\! = \bm{I}_d}_{\forall j\in [1, m]}\bigg\},
    \end{aligned}
\end{equation}

Here, for ease of notation, $\tilde{\bm{B}}^b(d,2^{m-i+1})$ is abbreviated as $\tilde{\bm{B}}^b_{(d,i)}$, and $\bm{I}_d$ denotes the $d$-dimensional identity matrix. The matrix $\bm{R}(m,b) \in \mathbb{R}^{d\times d}$ is an orthogonal matrix structured as a sequence of $m$ orthogonal butterfly matrices. For notational simplicity, we define BOFT instantiated with $\bm{R}(m,\frac{b}{2})$ as BOFT($m$, $b$) for $b\geq 2$. In the special case when $\thickmuskip=2mu \medmuskip=2mu m=1$, BOFT($1$, $b$) simplifies to the block-diagonal OFT~\cite{qiu2023controlling} using blocks of size $b$. When the block size $b=d$, BOFT($1$, $d$) becomes equivalent to the classical OFT~\cite{qiu2023controlling} utilizing an unrestricted orthogonal matrix. Alternatively, BOFT($\log \frac{2d}{b}$, $b$) utilizes the block butterfly matrix $\bm{B}^{\frac{b}{2}}(d)$ for $\bm{R}$, resulting in a dense orthogonal structure. More generally, BOFT($m$, $b$) comprises $\frac{1}{2}(b-1)dm$ trainable parameters when fine-tuning a linear transformation of dimensions $d \times n$. With the butterfly configuration ($m = \log d$, $b = 2$), BOFT requires only $\mathcal{O}(d\log d)$ parameters. Conversely, the standard OFT involving a dense full-rank orthogonal matrix entails $\mathcal{O}(d^2)$ parameters, and the block-diagonal OFT with $r$ blocks needs $\mathcal{O}(bd)$. Notably, to realize a dense orthogonal structure, OFT mandates $b=d$, while BOFT achieves this flexibility with any chosen $b$.

\textbf{Identity Initialization Strategy for BOFT.} Fine-tuning approaches generally initiate from the exact pretrained model, ensuring minimal deviation during adaptation; for instance, LoRA applies zero initialization to its low-rank component. In BOFT, we opt for initializing every butterfly factor as an identity matrix (i.e., setting the skew-symmetric matrices in the Cayley parametrization to zero).

\textbf{BOFT-specific Multiplicative Dropout.} Building upon LoRA~\cite{hulora2022}, which incorporates Dropout for the low-rank adaptation to avoid overfitting, the standard Dropout~\cite{srivastava2014dropout} mechanism fits well for LoRA but is incompatible with the multiplicative scheme in BOFT. To overcome this, we introduce a multiplicative Dropout tailored for BOFT. Because the orthogonal matrix $\bm{R}(m,b)$ consists of $m$ orthogonal butterfly components—each decomposable into $2b\times 2b$ block-diagonal orthogonal structures—we apply this Dropout by randomly masking $p_1$ percent of the butterfly components and $p_2$ percent of the blocks within each component, replacing the selected blocks with identity matrices.

\section{Intriguing Insights and Discussions}

\textbf{Expressivity of BOFT}. The butterfly structure, when combined with permutations, is capable of exactly representing a variety of well-known fast linear transforms~\cite{dao2019learning,dao2019kaleidoscope} (\emph{e.g.}, fast Fourier transform, Hadamard transform). Nevertheless, the extent to which an orthogonal butterfly matrix can approximate an arbitrary orthogonal matrix remains an open question. To investigate this, we carry out an experiment where a randomly generated dense orthogonal matrix~\cite{anderson1987generation} of dimension $512\times 512$ is approximated. The outcomes shown in Figure~\ref{fig:para_eff} reflect the mean across 10 independent trials with different random seeds. On the y-axis, we report the approximation error, while the x-axis presents the effective number of trainable parameters. Each color represents BOFT models sharing the same block size, and the leftmost point on each curve indicates the performance of BOFT($1$,$b$) (\emph{i.e.}, the standard block-diagonal OFT with block size $b$). BOFT consistently demonstrates superior parameter efficiency in comparison to OFT. For instance, BOFT($9$,$2$) achieves higher expressivity than BOFT($1$,$16$), despite using significantly fewer parameters. Typically, configurations with smaller $b$ and larger $m$ in BOFT show greater parameter efficiency. As another example, BOFT($6$,$4$) utilizes many fewer parameters while reaching a comparable approximation error to BOFT($2$,$16$). In summary, the butterfly matrix constitutes a more highly structured subset within the orthogonal group than block-diagonal architectures, establishing that BOFT is inherently more expressive than OFT given the same block size.

\begin{theorem}[Expressivity of BOFT]\label{thm:exp_boft}
    BOFT surpasses OFT in expressivity when using equivalent block sizes. To approximate any orthogonal matrix of dimension $d$, one can use a product of butterfly matrices of the form $\bm{B}_{d-1,1}(d)\bm{B}^\top_{d-1,2}(d)\cdots\bm{B}_{1,1}(d)\bm{B}^\top_{1,2}(d)$, where each $\bm{B}_{i,j}(d)$ represents a butterfly matrix.
\end{theorem}

Theorem~\ref{thm:exp_boft} directly implies a broader formulation for BOFT—specifically, the target orthogonal matrix can be generalized as $\thickmuskip=2mu \medmuskip=2mu \bm{R}^G(m_1,b_1,m_2,b_2,l)=\bm{R}_{l,1}(m_1,b_1)\bm{R}_{l,2}^T(m_2,b_2)\cdots\bm{R}_{1,1}(m_1,b_1)\bm{R}_{1,2}^T(m_2,b_2)$, where each $\bm{R}_{i,j}^T(m,b)$ indicates an orthogonal matrix utilized in BOFT. Setting $m_1=m_2=\log d$, $b_1=b_2=2$, and $l=d-1$, ensures that $\bm{R}^G(m_1,b_1,m_2,b_2,l)$ spans the full orthogonal group. This composition is referred to as a kaleidoscope hierarchy~\cite{dao2019kaleidoscope}. However, it is important to observe that increasing expressiveness does not always guarantee improved finetuning outcomes; for instance, full finetuning, although theoretically maximally expressive, may still perform suboptimally in practice. Thus, achieving a balance between expressivity and regularity is vital for ensuring model generalizability during finetuning. BOFT, by introducing structural priors, broadens the space of tunable parameters, thereby enabling a more favorable compromise between these two competing objectives.

\textbf{Spectral properties}. Orthogonal finetuning tends to maintain superior spectral characteristics compared to LoRA, primarily because it exactly preserves the spectral norm of the pretrained matrix $\bm{W}^0$. To illustrate this, consider the singular value decomposition $\bm{W}^0 = \bm{U} \bm{\Sigma} \bm{V}^\top$, where $\bm{U}$ and $\bm{V}$ are orthogonal matrices, and $\bm{\Sigma}$ is diagonal containing the singular values. Both OFT and BOFT apply an orthogonal transformation $\bm{R}$ from the left, resulting in finetuned weights of the form $\bm{R} \bm{U} \bm{\Sigma} \bm{V}^\top$. This procedure leaves the largest singular value (i.e., the spectral norm of $\bm{W}^0$) unchanged. Such preservation of the spectral norm has been demonstrated to improve both training stability and generalization performance~\cite{miyato2018spectral,yoshida2017spectral}. Further relevant mathematical details can be found in Appendix~\ref{app:sn_property}.

\textbf{Orthogonal finetuning as learning bilinear similarity}. The BOFT framework can be interpreted as optimizing for a bilinear similarity, specifically $\bm{w}^0_i\bm{R}\bm{x}$, where $\bm{w}^0_i$ denotes the $i$-th column (neuron) of $\bm{W}^0$. Here, BOFT's core is to learn a bilinear similarity matrix $\bm{R}$, subject to the strict requirement that $\bm{R}$ remain orthogonal. This connects directly to the literature on distance metric learning~\cite{xing2002distance} and the mathematical concept of a bilinear form~\cite{roman2005advanced}.

\textbf{Inductive bias and generalization in BOFT}. The matrix $\bm{R}(m,b)$ in BOFT frequently corresponds to a highly-structured subset of the orthogonal group, which, in effect, imposes a restriction on the range of possible models. Consequently, BOFT inherently generates an inductive bias. We posit that the particular inductive bias imposed through butterfly factorization is advantageous for generalization, owing to its alignment with common structures in various canonical linear transforms~\cite{dao2019kaleidoscope}, including the discrete Fourier, discrete sine/cosine, and Hadamard transforms. In addition, the sparse matrix decomposition characteristic of BOFT can introduce implicit forms of inductive bias~\cite{gunasekar2017implicit,li2018algorithmic,liu2019neural}.

\textbf{Comparison to butterfly-based sparse training}. Prior studies~\cite{dao2019learning,dao2019kaleidoscope,chen2022pixelated,dao2022monarch} have extensively explored sparse training through the butterfly parameterization approach. These works primarily reparameterize the weight matrices directly using the butterfly format and initiate neural network training from scratch. In~\cite{dao2022monarch}, a strategy for finetuning involves initially projecting pretrained weights onto a variant of butterfly matrices, after which the projected components are further optimized for downstream tasks. Distinct from these approaches, BOFT introduces a novel finetuning paradigm that leverages layer-shared weight matrices to transform the model parameters.

\input{rebuttal-results}

\section{Concluding Remarks and Limitations}

This work introduces Orthogonal Butterfly, a broadly applicable and parameter-efficient finetuning strategy for foundation models rooted in the butterfly architecture. The central concept for enhancing parameter efficiency lies in representing a dense orthogonal matrix as a product of multiple sparse orthogonal matrices. To systematically construct such matrix factorizations, we propose a framework based on graphical information transmission. Within this perspective, the butterfly configuration emerges as an effective structure for achieving our goals of sparse orthogonal matrix factorization. BOFT’s empirical performance is validated across several domains, including finetuning for large language models, large-scale vision models, and text-to-image generation systems. Additionally, our results consistently highlight BOFT's advantages as a general-purpose finetuning approach.

However, BOFT also has certain shortcomings. The computation during training is somewhat increased compared to OFT, since the resultant orthogonal matrix in BOFT is formed by multiplying several orthogonal matrices together. Improving BOFT's training efficiency thus remains an open challenge. Nonetheless, after finetuning, the learned orthogonal matrices can be fused into the original pretrained model, incurring no further latency during inference. Furthermore, it is still unclear if the butterfly architecture represents the optimal mechanism for information propagation. Our information transmission framework also allows us to draw insights from the domain of computer networking, where the design of efficient network topologies for information flow is a well-studied problem. We anticipate that adopting even more efficient network structures could further enhance the composition of dense orthogonal matrices.